
% 来源：EPUB part188.xhtml
\section{做对的事情，把事情做对}

\BookQuestion{网友：时常在想是什么原因让您跟我们这么不一样？本分和平常心。（2019-08-21）}

% 来源：EPUB part189.xhtml
\subsection{找到自己的北斗星}

做对的事情，大道最早（大三时）看到的是管理大师彼得·德鲁克说的。这是大道一生受益最大的一句话。（2021-08-12）

也许应该再看看蒂姆·库克2015年给华盛顿大学毕业生的讲话，努力去找到你自己的北斗星。我还想起巴菲特说过的，明知是坑就别再往下挖了。千万别拿时间不当成本。我见过太多聪明人忙碌一辈子也忙不明白的，原因大概就是没有自己的北斗星，凡事决策时都是从眼前利益着想，30年后的差异是巨大的。（2016-05-31）

北斗星指的是价值观，就是要有不为清单。30年后可以看到巨大的差异。（2016-06-01）

我没有过所谓的人生导师，但有颗北斗星一直在照着我哈。（2019-05-30）

曾经有个博友名字叫“守正出奇”，我说我不喜欢这个名字。问为什么？我说守正出奇的人其实都是整天在想着出奇的。网友问怎么改？我说：守正不出奇。

勤奋和天赋其实都没那么重要，做对的事情最重要。（2019-09-09）

\BookQuestion{网友：如果勤奋和天赋都不够，那怎么搞明白什么是对的事？}

如果我说勤奋重要你会变得更勤奋吗？如果我说天赋重要你会变得更有天赋？但是，做对的事情是可以选择的。人们大概都知道骗人是不对的，也大概知道抽烟是不对的吧，也大概知道酗酒是不对的吧？看看周围有多少人在选择做对的事情你大概就明白了。（2019-09-10）

“我们长期努力保持不做傻事，所以我们的收获比那些努力做聪明事的人要多得多。”（芒格）

小聪明和大智慧的差别，也是坚持不做不对的事情的意思。

“聪明人也不免遭受过度自信带来的灾难。他们认为自己有更强的能力和更好的方法，所以往往他们就在更艰难的道路上疲于奔命。”（芒格）

最简单的例子就是一个很会开车的人开了一辆很好的车但上高速时上反了入口哈。

指的是聪明反被聪明误的那种？其实就是做对的事情（坚持不做错的事情），然后把事情做对。聪明往往指的是把事情做对的能力比较强，但知道坚持不做错的事情的人最后才会成大器，比如巴菲特。

大家一定要分清做错的事情和在把事情做对的过程中犯错误的本质区别。在把事情做对的过程当中任何人都是会犯错的，包括巴菲特，但由于坚持不做错的事情，所以犯大错的概率低，很多年以后的结果就完全不一样了。如果还不明白的话，就请看看那些一直很聪明但老是不如意的人们这些年都在忙啥就明白了。（2012-06-26）

往回看个几十年，你会看到很多很“聪明”的“聪明人”成就很小，原因可能在他们没把聪明放在做对的事情上。（2015-04-06）

% 来源：EPUB part190.xhtml
\subsection{想长远，想本质}

\BookQuestion{网友：大道，您离开北京的国企、离开小霸王、去美国陪夫人，每一个决定换作别人都会纠结很久，但您能很理性地处理沉没成本，这种能力是经过多少时间形成的，有没有突然“开悟”的诱因？}

也许是一种习惯，就是想长远的习惯。不知道怎么开始的了，回想起来也许是本科毕业考研时，找了3天找不到自己未来想做的事情（找不到想考的专业），就把考研放弃了。（2020-12-03）

\BookQuestion{网友：在离开北京电子管厂、佛山无线电厂的两次工作转变中，从发现不对到做出改变用了多久？}

想清楚不对就马上决定改变了，不过改变是需要时间的。比如考研需要备考，换工作需要找工作，等等。（2020-12-04）

大约大三时悟到我大概不适合当个科学家或者工程师啥的，同学里有些人读书就是比我厉害。放眼全世界，那得有多少人比我厉害啊，但我知道的科学家屈指可数，加上内心对所学的东西总是提不起兴趣（这也可能是最重要的原因），所以那时候就开始想我也许不会成为一个工程师。大学毕业时本来也为考研准备了很久，但到报名时居然找不到自己想考的专业，所以最后就放弃了。那几年真是觉得很迷惘，但冥冥中觉得不对的事情就不应该继续下去了。我可能是那个年代很少有的工科毕业但没干过一天技术工作的人。（2021-06-30）

\BookQuestion{网友：我今年30岁，感觉自己太急躁，最近几年事业又是低谷期。您遇到这种问题如何解决？}

最重要的是做任何选择要有是非感而不是跟着利益走。这么做的话， 过些年就会好很多的。（2019-03-15）

什么时候意识到要做对的事情都不晚。不过，满脑子只有利益没有对错的人们确实是很难从原来的坑里出来的。（2019-05-28）

\BookQuestion{网友：对理性的定义是什么？}

理性就是想长远啊。（2022-03-08）

\BookQuestion{网友：您的回答很快就能说出问题的本质。这种能力我们该如何学习？}

想学就可以，只要努力看远一点就好。

大多数人其实是不在乎问题的本质的。（2020-11-17）

\BookQuestion{网友：根据您的切身体会，什么才是提升悟性的关键？}

我觉得没人天生有悟性的。我看到的有悟性的人一般都非常肯悟，会习惯想本质及长远。（2020-10-23）

“习惯的枷锁，开始的时候轻得难以察觉，到后来却重得无法摆脱。”（巴菲特）

其实这条非常非常重要，不要等到太晚才看明白啊。（2012-11-19）

\BookQuestion{网友：商业直觉是天赋吗？}

所谓商业直觉其实就是理性思考能力，巴菲特早就是顶级的了。当然，随着经验的积累，对不同公司的直觉的准确度肯定是会增加的。比如我能比较懂手机，巴菲特比较懂保险公司和银行等等。巴菲特懂的东西太多，我这么说不是很全面哈。（2023-01-11）

\BookQuestion{网友：如果列出一条您成功的必然因素和一条偶然因素，请问会是什么？}

没有必然因素，偶然因素太多了。概率上来说，因为我习惯想长远，确实增加了很多成功的概率。（2020-12-03）

\BookQuestion{网友：我总觉得看问题看不到点子上，广泛撒网阅读，博取百家看法，最后还是很难得到确定性，看问题进不到下一层。}

如果你能用平常心去看问题就会容易看到事物的本质。（2011-03-28）

\BookQuestion{网友：你的平常心是怎么来的？}

平常心其实就是理性思考的习惯，应该是后天养成的。（2019-04-22）

\BookQuestion{网友：这两年，学习了在网上能查到的您的一切，精神和物质都有很大收获。我现在面对诱惑时，依然有时会失去平常心，但经常反省，继续努力。}

告诉你一个好消息，我之所以提出平常心就是因为我也会失去平常心，所以大家都差不多，需要时刻提醒一下自己哈。（2019-09-12）

平常人是很难有平常心的，所以平常心又叫不平常心。如果你关心的不是事物的本质，没有平常心是正常的。倒过来也一样，如果你关注的只是事物的本质，平常心自然就在那里了。在错的道路上是没办法有平常心的。（2015-05-29）
